\section{Development case studies}

\subsection{Case P40: propagation buckling in PIP systems}
The P40 case is based on the propagation-buckling study of Karampour et al. \cite{Karampour2017}. The automated study selected a source-parameter table, a comparison table, and the main PIP propagation-pressure relationship. The analytical method was reproduced directly, while an independent work-balance reconstruction provided a second mechanics path.

For two of the three PIP configurations, the reproduced normalized propagation-pressure ratio agrees with the published Table 2 value to rounding. For the third case, the reproduced ratio is approximately 0.7057 while the published value is 0.66. An independent reconstruction gives approximately 0.7053, corroborating the reproduced value rather than the printed table value. The disagreement remains open and is not removed through parameter adjustment.

\begin{table}[htbp]
\centering
\caption{P40 selected reproduction results. Values are retained as development evidence, not as a claim that the source publication is erroneous.}
\label{tab:p40}
\begin{tabular}{lccc}
\toprule
Case & Published Eq. (9) ratio & Reproduced ratio & Independent check \\
\midrule
PIP-1 & 0.86 & 0.8635 & consistent \\
PIP-2 & 0.69 & 0.6900 & consistent \\
PIP-3 & 0.66 & 0.7057 & $\approx 0.7053$ \\
\bottomrule
\end{tabular}
\end{table}

\subsection{Case P41: axial load sharing and global buckling}
The P41 case is based on Goplen et al. \cite{Goplen2011}. This source tests different mechanics and a different publisher format. The first evidence chain uses the published pipe geometry and Young's moduli to independently reconstruct steel area and axial stiffness, then applies the published Eqs. (6)--(8) to the soil-friction load split.

The independently calculated steel areas and axial stiffnesses reproduce the published Table 2 values closely. For a published soil-friction increment of 471~N/m, the load-sharing equations give approximately 158.8~N/m for the inner pipe and 312.2~N/m for the outer pipe, closing to 471~N/m. The printed table instead reports 153 and 312~N/m.

A second P41 chain evaluates the published full-bonding friction criterion. With $\gamma_f=1$ for the paper's illustrative case classification, the required internal friction is approximately 158.78~N/m. This independently reproduces the four published Table 3 categories: no friction, partial bonding, full bonding, and full bonding. The paper's dry-friction estimate is also independently checked: $1195\times0.3=358.5$~N/m, consistent with the rounded 358~N/m value used in the highest-friction case.

The final P41 checkpoint reviews the global-buckling response and Eq. (14). An elastic geometry check gives
\[
\frac{(EI)_{\mathrm{outer}}}{(EI)_{\mathrm{inner}}}\approx 3.3966,
\]
which independently supports the paper's qualitative statement that the outer pipe is normally the stiffer bending member. This check does not reproduce the nonlinear moment--rotation curves or the finite-element model.
